\documentclass[10pt]{article}
\usepackage[margin=0.75in]{geometry}
\usepackage{amsmath,amssymb,amsthm,microtype,hyperref}
\newtheorem{theorem}{Theorem}
\newcommand{\B}{\mathcal B}
\newcommand{\R}{\mathcal R}
\begin{document}
\begin{center}
{\Large Exact Affine-Transfer Reformulations of Collatz}\par\medskip
Collatz proof project\quad---\quad September 5, 2026
\end{center}
Let $T(n)=n/2$ for even $n$ and $T(n)=(3n+1)/2$ for odd $n$.
Write $\B(n)$ for boundedness of its natural orbit and $\R(n)$ for
the existence of $t\geq0$ with $T^t(n)=1$. The previously proposed
transfer $z\mapsto9z+2$ is not established here. We determine its exact
logical strength, including the effect of restricting its domain.

\begin{theorem}[Nondivergence]
The following three assertions are equivalent:
\begin{align*}
&\forall n\geq0,\ \B(n);\\
&\forall z\geq0,\ \B(z)\Longrightarrow\B(9z+2);\\
&\forall z\equiv2\pmod3,\ \B(z)\Longrightarrow\B(9z+2).
\end{align*}
Moreover, any unbounded natural orbit implies the existence of a
$z\equiv2\pmod3$ with $\B(z)$ and $\neg\B(9z+2)$.
\end{theorem}

\begin{theorem}[Full conjecture]
The Collatz conjecture is equivalent to restricted reaching-one transfer:
\[
\bigl(\forall n>0,\ \R(n)\bigr)
\quad\Longleftrightarrow\quad
\bigl(\forall z\equiv2\pmod3,\ \R(z)\Longrightarrow\R(9z+2)\bigr).
\]
This version includes the exclusion of nontrivial cycles. Neither side
is proved by the equivalence.
\end{theorem}

\begin{proof}[Common argument]
Use either a least unbounded seed, or strong induction on a seed whose
convergence is to be established. Zero is bounded; one reaches one.
Every positive even seed has a smaller positive successor. For an odd
$n>1$, write $n+1=2^am$ with $m$ odd. If $a=1$, two steps give
$(3n+1)/4<n$. Thus only $a\geq2$ remains.

The classical odd-run identities, formalized in \texttt{OddRunMerges}, give
two cases at time $a+2$. If $3^am\equiv1\pmod4$, the seeds $n,n-1$
merge. The smaller seed supplies boundedness or convergence of the
common tail. Otherwise, putting $z=T^{a+2}(n-1)$ gives
\[
 T^{a+2}(n)=9z+2,\qquad 4z+1=3^{a-1}m.
\]
Since $a\geq2$, the second identity gives $z\equiv2\pmod3$.
The positive seed $n-1$ supplies $\B(z)$ or $\R(z)$; the restricted
transfer supplies the corresponding property of $T^{a+2}(n)$ and
hence of $n$.

For boundedness, if any orbit were unbounded, its least seed would
therefore produce the asserted failure pair $z,9z+2$. For convergence,
this reasoning is a strong-induction proof using the transfer as a
hypothesis. Reaching one is invariant under finite shifts, including
shifts past an earlier hit of one: the orbit then stays in the cycle
$1,2$. Conversely, universal boundedness or universal convergence
immediately implies its corresponding transfer statement.
\end{proof}

\paragraph{Formal evidence.}
\texttt{Collatz.AffineBoundedness} defines \texttt{OrbitBounded}, constructs
the bounded/unbounded failure pair, and proves both boundedness equivalences.
\texttt{Collatz.AffineConvergence} defines \texttt{ReachesOne}, proves shift
invariance, and proves the full-conjecture equivalence by strong induction.
The supporting residue theorem is in \texttt{Collatz.OddRunMerges}.
All eight added or newly selected declarations are included in the axiom
audit. No numerical search supplies an unproved premise.

\paragraph{Interpretation.}
Restricting transfer to $z=2$ modulo three does not make it a consequence
of the existing local identities. It still carries the entire target
statement. In the reaching-one version the equivalent form is
$\R(3u+2)\Rightarrow\R(27u+20)$ for every natural $u$.
A proof must establish that implication; testing many convergent pairs
or merely observing an affine relationship is insufficient. This note
is an exact reformulation and a check against circular reasoning, not
a proof of Collatz or a claim of mathematical priority.

\paragraph{Reproduction and background.}
\begingroup\raggedright
From \texttt{lean/}, build \texttt{Collatz.AffineBoundedness} and
\texttt{Collatz.AffineConvergence}; from the root run the theorem audit
with \texttt{--build}. The classical endpoint identities used above are
also in M. D. LaDue, \emph{Clusters of Integers with Equal Total Stopping
Times in the $3x+1$ Problem} (2017), Lemma 3.1 and Theorem 4.1,
\href{https://arxiv.org/abs/1709.02979}{arXiv:1709.02979}.\par\endgroup
\end{document}
